% n-Queens Completion as the energy of a {0,1} Boltzmann machine, certified in Lean 4, with the
% collaborative neurodynamic algorithms of Li and Wang run on it.  Everything here is
% HopfieldNet/QUBO/ and HopfieldNet/QUBO/Instances/Queens/.
%
% Written in the format and register of a short conference paper: rules, translation, what is
% proved, the search, results, concluding remarks.  Nothing is claimed that the development does not
% carry; where the certification stops is stated rather than elided.
%
% Provenance of the numbers.  The two tables come from
%   lake build HopfieldNet.QUBO.Instances.Queens.Bench   -- ~16 min
% and the encoding counts are cross-checked independently by a script under scripts/.
%
% Every board drawn below is one of the corpus instances of
%   HopfieldNet/QUBO/Instances/Queens/Bench.lean:80
% and is restated as a #guard in Figures.lean.  The chessboards use the same normalised outline and
% palette as the interactive widget, so a figure here and a frame in the infoview are the same
% picture.
%
% NOTE: copy `easychair.cls` and `lstlean.tex` from the main paper's directory next to this file.
\documentclass{easychair}

\usepackage{amsmath}
\usepackage{amssymb}
\usepackage{booktabs}
\usepackage{caption}
% `H' pins a float exactly where it is written, instead of letting it drift to the top of a page.
\usepackage{float}
\usepackage{listings}
\usepackage{xcolor}
\usepackage{csquotes}
\usepackage{fontawesome}
\usepackage{tikz}

% The environment structure of the HNBM paper: theorems, lemmas, propositions, corollaries and
% remarks share one counter; definitions carry their own.
\newtheorem{thm}{Theorem}[section]
\newtheorem{theorem}[thm]{Theorem}
\newtheorem{lemma}[thm]{Lemma}
\newtheorem{proposition}[thm]{Proposition}
\newtheorem{corollary}[thm]{Corollary}
\newtheorem{remark}[thm]{Remark}
\newtheorem{definition}{Definition}[section]
% A `proof' environment without amsthm, so that nothing depends on the class's theorem package.
\makeatletter
\@ifundefined{proof}
  {\newenvironment{proof}{\par\noindent\textit{Proof.}\ }{\hfill$\square$\par\medskip}}{}
\makeatother

\definecolor{academicpurple}{rgb}{0.4,0.0,0.35}
\newcommand{\secref}[1]{\hyperref[#1]{\textcolor{academicpurple}{\S\ref*{#1}}}}
\newcommand{\citeref}[1]{\textcolor{academicpurple}{\cite{#1}}}
\newcommand{\thmref}[1]{\textcolor{academicpurple}{\hyperref[#1]{Theorem~\ref*{#1}}}}
\newcommand{\lemref}[1]{\textcolor{academicpurple}{\hyperref[#1]{Lemma~\ref*{#1}}}}
\newcommand{\propref}[1]{\textcolor{academicpurple}{\hyperref[#1]{Proposition~\ref*{#1}}}}
\newcommand{\corref}[1]{\textcolor{academicpurple}{\hyperref[#1]{Corollary~\ref*{#1}}}}
\newcommand{\remref}[1]{\textcolor{academicpurple}{\hyperref[#1]{Remark~\ref*{#1}}}}
\newcommand{\defref}[1]{\textcolor{academicpurple}{\hyperref[#1]{Definition~\ref*{#1}}}}
\newcommand{\tabref}[1]{\textcolor{academicpurple}{\hyperref[#1]{Table~\ref*{#1}}}}
\newcommand{\figref}[1]{\textcolor{academicpurple}{\hyperref[#1]{Figure~\ref*{#1}}}}
\newcommand{\extlink}{~\ensuremath{{}^\text{\faExternalLink}}}
%%% Source links.  `\reporef' is the commit the line numbers below were taken from; every anchor in
%%% this paper was checked against it.  For the links to resolve, that commit must be reachable
%%% from a ref pushed to `\repoowner' -- as of writing it is NOT: the queens development lives only
%%% on the local branch `bump-mathlib-4.32', and neither `origin/main' nor the fork contains
%%% `HopfieldNet/QUBO/Instances/Queens/'.  Push it, then run
%%%   bash scripts/check_paper_links.sh papers/lpar-queens/sat-queens.tex
%%% which fetches every link and fails on the first one that does not return 200.  If the commit
%%% changes, re-run that script: the line numbers move with it.
\newcommand{\repoowner}{mkaratarakis/HNBM}
\newcommand{\reporef}{1bd91612ee31180f6c99fa11090b713f3012ead4}
\newcommand{\repo}{https://github.com/\repoowner/blob/\reporef}
\newcommand{\src}[1]{\href{\repo/#1}{\extlink}}
% \srcl{path}{line} -- the clickable symbol, anchored at the declaration itself.
\newcommand{\srcl}[2]{\href{\repo/#1\#L#2}{\extlink}}
% Shorthands for the two directories every link below sits in.
\newcommand{\qsrc}[2]{\srcl{HopfieldNet/QUBO/Instances/Queens/#1}{#2}}
\newcommand{\qusrc}[2]{\srcl{HopfieldNet/QUBO/#1}{#2}}

\definecolor{keywordcolor}{rgb}{0.7, 0.1, 0.1}
\definecolor{tacticcolor}{rgb}{0.0, 0.1, 0.6}
\definecolor{commentcolor}{rgb}{0.4, 0.4, 0.4}
\definecolor{symbolcolor}{rgb}{0.0, 0.1, 0.6}
\definecolor{sortcolor}{rgb}{0.1, 0.5, 0.1}
\definecolor{attributecolor}{rgb}{0.7, 0.1, 0.1}
\definecolor{codebg}{rgb}{0.96,0.96,0.99}
\definecolor{codeframe}{rgb}{0.75,0.70,0.80}
\def\lstlanguagefiles{lstlean.tex}
\lstset{language=lean,
  basicstyle=\fontsize{8.4}{9.8}\selectfont\ttfamily,
  backgroundcolor=\color{codebg},
  frame=single, framerule=0.5pt, rulecolor=\color{codeframe},
  framesep=4pt, xleftmargin=6pt, xrightmargin=6pt,
  aboveskip=4pt, belowskip=4pt}

%%% ------------------------------------------------------------------- chessboards
%%% Palette and queen outline taken from queensBoard.js, so that the figures below and the widget
%%% are the same picture.  A square is addressed (column, -row) with row 0 at the top, matching the
%%% representation used throughout: q_i is the column of the queen in row i.
\definecolor{sqlight}{HTML}{F0D9B5}
\definecolor{sqdark}{HTML}{B58863}
\definecolor{sqedge}{HTML}{8A6242}
\definecolor{attackc}{HTML}{C0392B}
\definecolor{qdarkbody}{HTML}{26221E}
\definecolor{qdarkedge}{HTML}{0D0B0A}
\definecolor{qlitebody}{HTML}{FBFAF7}
\definecolor{qliteedge}{HTML}{4A423A}
\definecolor{cRow}{HTML}{1F5FA8}
\definecolor{cCol}{HTML}{1E7A48}
\definecolor{cDiag}{HTML}{B26B00}
\definecolor{cAnti}{HTML}{6B3FA0}

%%% Every macro below is prefixed `qc' so that nothing here can collide with the document class or
%%% with a package loaded later.  `qcpic' opens a picture whose baseline is its centre, so boards of
%%% different sizes line up when set side by side in a tabular.
\newenvironment{qcpic}[1]
  {\begin{tikzpicture}[x=#1,y=#1,baseline=(current bounding box.center)]}
  {\end{tikzpicture}}

% \qcboard{N} -- the squares and the frame of an N x N board.
\newcommand{\qcboard}[1]{%
  \pgfmathtruncatemacro{\qcn}{#1-1}%
  \foreach \qci in {0,...,\qcn}{\foreach \qcj in {0,...,\qcn}{%
    \pgfmathtruncatemacro{\qcmix}{mod(\qci+\qcj,2)==0 ? 100 : 0}%
    \fill[sqlight!\qcmix!sqdark] (\qcj,-\qci) rectangle ++(1,-1);}}%
  \draw[sqedge,line width=0.5pt] (0,0) rectangle (#1,-#1);}

% \qcglyph{col}{row}{body}{edge} -- one queen, from the widget's outline.
\newcommand{\qcglyph}[4]{%
  \begin{scope}[shift={(#1,-#2)},line join=round,line width=0.35pt,draw=#4,fill=#3]
    \filldraw (0.14,-0.255) -- (0.265,-0.45) -- (0.32,-0.185) -- (0.41,-0.45)
      -- (0.50,-0.155) -- (0.59,-0.45) -- (0.68,-0.185) -- (0.735,-0.45)
      -- (0.86,-0.255) -- (0.80,-0.555) -- (0.20,-0.555) -- cycle;
    \filldraw[rounded corners=0.5pt] (0.175,-0.555) rectangle (0.825,-0.640);
    \filldraw (0.255,-0.64) .. controls (0.255,-0.78) and (0.205,-0.83) .. (0.165,-0.875)
      -- (0.835,-0.875) .. controls (0.795,-0.83) and (0.745,-0.78) .. (0.745,-0.64) -- cycle;
    \filldraw[rounded corners=0.5pt] (0.125,-0.875) rectangle (0.875,-0.950);
    \foreach \qcu/\qcv in {0.14/0.235, 0.32/0.165, 0.50/0.135, 0.68/0.165, 0.86/0.235}
      \filldraw (\qcu,-\qcv) circle (0.072);
  \end{scope}}
% A given queen is dark, a queen the solver placed is light -- the widget's convention.
\newcommand{\qcdark}[2]{\qcglyph{#2}{#1}{qdarkbody}{qdarkedge}}
\newcommand{\qclite}[2]{\qcglyph{#2}{#1}{qlitebody}{qliteedge}}
% \qcgivens{i/j,...} -- reads exactly like the Lean array literal #[(i,j),...].
\newcommand{\qcgivens}[1]{\foreach \qcgi/\qcgj in {#1}{\qcdark{\qcgi}{\qcgj}}}
% \qcplace{q_0,...,q_{N-1}} -- reads exactly like the Lean array #[q_0,...,q_{N-1}].
\newcommand{\qcplace}[1]{%
  \foreach \qccol [count=\qcrow from 0] in {#1}{\qclite{\qcrow}{\qccol}}}
% \qcrule{colour}{i}{j}{i'}{j'} -- one row of A-hat, as a capsule through the centres of its squares.
\newcommand{\qcrule}[5]{%
  \draw[#1,line cap=round,opacity=0.42,line width=0.30cm]
    ({#3+0.5},{-#2-0.5}) -- ({#5+0.5},{-#4-0.5});}
% \qcmark{i}{j} -- outline one square.
\newcommand{\qcmark}[2]{\draw[black,line width=0.9pt] (#2,-#1) rectangle ++(1,-1);}
% \qcattack{i}{j}{i'}{j'} -- the widget's attack indicator: tinted squares and a dashed segment.
\newcommand{\qcattack}[4]{%
  \fill[attackc,opacity=0.28] (#2,-#1) rectangle ++(1,-1);
  \fill[attackc,opacity=0.28] (#4,-#3) rectangle ++(1,-1);
  \draw[attackc,line width=0.8pt,dash pattern=on 2.4pt off 1.4pt]
    ({#2+0.5},{-#1-0.5}) -- ({#4+0.5},{-#3-0.5});}
% \qcnames{N} -- the name of the variable of every square, in the square.
\newcommand{\qcnames}[1]{%
  \pgfmathtruncatemacro{\qcn}{#1-1}%
  \foreach \qci in {0,...,\qcn}{\foreach \qcj in {0,...,\qcn}{%
    \node[font=\tiny] at ({\qcj+0.5},{-\qci-0.5}) {$x_{\qci\qcj}$};}}}
\newcommand{\qclab}[1]{{\footnotesize #1}}

% Loaded last, and only if the class has not already done so: \secref and \src need \hyperref and
% \href, and an undefined one of those is a hard error rather than a bad box.
\usepackage{hyperref}

\title{A Certified Neurodynamic Solver for $n$-Queens Completion}

\titlerunning{A Certified Neurodynamic Solver for $n$-Queens Completion}

\author{
Matteo Cipollina\inst{1}
\and
Michail Karatarakis\inst{2}
\and
Freek Wiedijk\inst{2}
}

\institute{
  Axiomatic AI, Barcelona, Spain\\
  \email{matteo.cipollina@yahoo.it}
\and
  Radboud University, Nijmegen, The Netherlands\\
  \email{michail.karatarakis@ru.nl, freek@cs.ru.nl}\\
}

\authorrunning{Cipollina, Karatarakis and Wiedijk}

\begin{document}
\maketitle

\begin{abstract}
This paper presents a certified neurodynamic solver for $n$-Queens Completion. An instance is
translated into a quadratic objective on $N^{2} + 4N - 6$ binary neurons whose minimum is attained
if and only if the instance has a completion, and a board can be read off any state attaining it.
The objective is proved to be the energy of a $\{0,1\}$ Boltzmann machine, and every infeasible
state is proved to lie at least $\tfrac12$ above the minimum, so a population of such machines run
collaboratively --- with a particle swarm rule re-initializing them whenever the population
converges --- applies to it unchanged. Little coding was necessary to implement this solver: the
network, its dynamics and the collaborative search are supplied by an existing formalization of the
two models in the Lean~4 proof assistant, and only the constraints on a completion had to be
specified in the prover's logic. The translation itself, however, is not left to inspection. It is
proved sound and complete, and the Boolean checker against which those properties are stated is in
turn proved to decide the chess condition written over the integers, so a run that reaches the
minimum establishes a theorem about the board. Results on eleven completable and seven blocked
instances are reported against a certified backtracking baseline; on six of the seven blocked
instances the absence of a completion is itself a theorem, and what the certification does not
cover is stated rather than elided.
\end{abstract}

\section{Introduction}\label{sec:intro}

$n$-Queens Completion is a placement puzzle. Given an $N \times N$ chessboard --- most frequently
the ordinary $8 \times 8$ board --- with queens standing on some of its squares (the
\enquote{givens}), the aim is to place further queens until there are $N$ of them, so that no two
share a row, a column or a diagonal. \figref{fig:example} shows an instance on the left, along with
its unique completion on the right. Note that the columns of the board, unlike the digits of a
number puzzle, cannot be replaced by arbitrary symbols: two queens share a diagonal exactly when
their row and column indices have equal difference or equal sum, so the arithmetic of those indices
is part of the puzzle. The plain problem, in which nothing is given, has been studied since 1848 and has a
solution for every $N \ge 4$, by explicit construction~\citeref{bell2009}.

The completion problem is harder than the plain one, and harder than it looks. Deciding whether an
instance has a completion is \textsf{NP}-complete, and counting the completions is
$\#\textsf{P}$-complete~\citeref{gent2017}. Whether a completion exists depends on the givens in a
way that is not local: two queens that share no row, no column and no diagonal can already rule one
out, and \figref{fig:hard} shows a $7 \times 7$ board on which they do. Solvers for the problem
exist, and at the sizes considered here backtracking with constraint propagation settles an instance
quickly. What it does not settle is the negative answer. A search that finds no completion has
established nothing about the board, and for this problem the instances worth posing are exactly
those that have none.

In this paper we propose a neurodynamic approach. Discrete Hopfield networks~\citeref{hopfield1982}
and Boltzmann machines~\citeref{ackley1985} settle into states of low energy, so a constraint problem
can be attacked by writing its constraints as an energy and letting a network descend it. An
instance of $n$-Queens Completion is translated into such an energy, on $N^{2} + 4N - 6$ binary
neurons, whose minimum is attained if and only if the instance has a completion; a board can then be
read off the state the network stops at. A single network stops at the first state no update
improves, which is rarely a minimum, so we use the collaborative scheme of Li and
Wang~\citeref{liwang2022}: a population of networks is run concurrently and their initial states are
re-drawn by a particle swarm rule whenever the population converges.

The translation is simple and required little implementation effort, since the energy, the network
it defines, its dynamics and the collaborative search are supplied unchanged by an existing
formalization of the two models in the Lean~4 proof assistant~\citeref{lean4,hnbm}; only the
constraints of the puzzle had to be specified. What did take effort is that the translation is not
left to inspection. It is proved sound and complete, the checker against which those properties are
stated is itself proved to decide the chess condition written over the integers, and the objective
is proved to be the energy of the network said to be descending it. A run that ends at the minimum
therefore establishes a theorem about the board; a run that ends above it establishes nothing, and
the boards with no completion are refuted by other means, on six of the seven in \tabref{tab:blocked}
by proof.

\begin{figure}[H]
\centering
\begin{tabular}{@{}c@{\hspace{2.6em}}c@{}}
\begin{qcpic}{0.46cm}
  \qcboard{8}\qcgivens{0/0,1/4,2/7}
\end{qcpic}
&
\begin{qcpic}{0.46cm}
  \qcboard{8}\qcplace{0,4,7,5,2,6,1,3}\qcgivens{0/0,1/4,2/7}
\end{qcpic}
\\[0.5em]
\qclab{(a) an instance} & \qclab{(b) its completion}
\end{tabular}
\caption{An $n$-Queens Completion instance and its solution. Given queens are dark and queens the
solver placed are light, which is the convention of the development's board widget. This instance
has exactly one completion.}\label{fig:example}
\end{figure}

\section{Implementation in Lean 4}\label{sec:rules}

An implementation of the rules of $n$-Queens Completion in the Lean~4 proof assistant is
straightforward. An \emph{instance} is a pair $(N, G)$ with $N \in \mathbb{N}$ the side of the
board and $G \subseteq \{0,\dots,N-1\}^{2}$ the givens, listed as (row, column) pairs. A board is
modelled by an array $q$ of natural numbers, $q_i$ denoting its entry at $i$ --- the column of the
queen in row $i$ --- read with the out-of-range default $N$, so that a missing entry violates the
bound in \defref{def:completion}(ii) below rather than passing silently. One queen per row is thus
built into the representation, and what remains is a condition on pairs of rows. We say that two
queens \emph{attack} each other iff they stand in the same column, on the same difference diagonal,
or on the same sum diagonal.

\begin{definition}[attack]\label{def:attack}
For $i, a, k, b \in \mathbb{Z}$ (\lstinline{Attack})\qsrc{Spec.lean}{37},
\[
  \mathrm{Attack}(i,a,k,b)
  \quad:\Longleftrightarrow\quad
  a = b \;\;\lor\;\; i - a = k - b \;\;\lor\;\; i + a = k + b .
\]
\end{definition}

Indexing the $N^{2}$ squares of the board as in \figref{fig:vars}, we can now say what it means for
$q$ to complete the instance: it has one entry per row, every entry is a column of the board, no
two queens attack, and every given is present.

\begin{definition}[completion]\label{def:completion}
An array $q$ of natural numbers is a \emph{completion} of $(N,G)$
(\lstinline{IsCompletion})\qsrc{Spec.lean}{44} iff
\begin{enumerate}
\item[(i)]   $\lvert q\rvert = N$;
\item[(ii)]  $q_i < N$ for every $i < N$;
\item[(iii)] $\lnot\,\mathrm{Attack}(i, q_i, k, q_k)$ for all $i, k < N$ with $i \ne k$, the four
             indices being taken in $\mathbb{Z}$;
\item[(iv)]  $q_r = c$ for every $(r,c) \in G$.
\end{enumerate}
\end{definition}

\defref{def:attack} is stated over $\mathbb{Z}$, so that its two diagonal conditions are literal
differences and sums; in (iii) the natural-number indices are coerced. That is not how the
development computes, and the discrepancy is the reason the next theorem is worth proving.

\begin{theorem}[the checker decides the condition]\label{thm:spec}
For every instance and every array $q$, the executable checker
\lstinline{isQueens}\qsrc{Basic.lean}{613} returns \lstinline{true} if and only if $q$ is a completion
in the sense of \defref{def:completion} (\lstinline{isQueens_iff})\qsrc{Spec.lean}{57}.
\end{theorem}

\begin{remark}\label{rem:checker}
Two features of the checker make \thmref{thm:spec} more than bookkeeping. First, it indexes squares
by natural numbers, on which subtraction is truncated --- $2 - 5$ is $0$, not $-3$ --- so the
literal transcription of $i - a = k - b$ is wrong, and the checker tests the rearranged form
$i + q_k \ne k + q_i$ instead. That rearrangement is valid over $\mathbb{Z}$ and is exactly the kind
of step that is easy to get backwards. Second, clauses (ii) and (iv) are stated through the
defaulting accessor described above, which is a coding convenience and not a statement about chess.
After \thmref{thm:spec} the specification a reader has to trust is \defref{def:completion}, which
mentions no arrays, no defaults and no truncated arithmetic. As an independent check the checker is
compared against a separately written implementation over $\mathbb{Z}$
(\lstinline{naive}\qsrc{Spec.lean}{130}) --- a list-based pairwise scan with real subtraction, no array
and no default --- exhaustively on all $4^{4}$ boards of size four and all $5^{5}$ of size five
(\lstinline{isQueens_eq_naive_four}\qsrc{Spec.lean}{141},
\lstinline{isQueens_eq_naive_five}\qsrc{Spec.lean}{145}).
\end{remark}

The next section translates \defref{def:completion} into an energy.

\begin{figure}[H]
\centering
\begin{tabular}{@{}c@{\hspace{2.6em}}c@{}}
\begin{qcpic}{0.86cm}
  \qcboard{5}\qcnames{5}
\end{qcpic}
&
\begin{qcpic}{0.86cm}
  \qcboard{5}
  \qcrule{cRow}{2}{0}{2}{4}
  \qcnames{5}
\end{qcpic}
\\[0.5em]
\qclab{(a) the squares} & \qclab{(b) one board row}
\end{tabular}
\caption{The squares of a $5 \times 5$ board and their neurons. (a) $x_{ij} = 1$ means that a queen
stands in row $i$ and column $j$. (b) The five neurons of board row $2$ carry one queen between
them, which is the constraint~\eqref{eq:form-row}; the capsule is the row of $A$ that says
so.}\label{fig:vars}
\end{figure}

\section{Translation to a Boltzmann machine}\label{sec:qubo}

Fix an instance $(N,G)$ with $N \ge 2$. We encode a board by introducing one binary neuron for each
square of the grid, i.e.\ $N^{2}$ neurons, together with one auxiliary neuron for each diagonal that
can carry a conflict. Each neuron $x_{ij}$ (with $0 \le i, j < N$) represents the truth value of the
statement that a queen stands on the square in row $i$ and column $j$; each auxiliary neuron records
that its diagonal is empty. For $0 \le d \le 2N-4$ put
\begin{equation}\label{eq:diags}
  D_d = \{(i,j) : 0 \le i,j < N,\ i - j = d + 2 - N\},
  \qquad
  A_d = \{(i,j) : 0 \le i,j < N,\ i + j = d + 1\} .
\end{equation}
Since $i - j$ ranges over $-(N-1),\dots,N-1$ and $i + j$ over $0,\dots,2N-2$, the two families
\eqref{eq:diags} enumerate exactly the diagonals containing at least two squares: the four omitted
values are the extremes, and each of those diagonals is a single corner, which cannot carry two
queens. There are $2N-3$ of each. The rules of \defref{def:completion} then become the linear system
\begin{subequations}\label{eq:form}
\begin{align}
  \textstyle\sum_{j} x_{ij} &= 1, & i &= 0,\dots,N-1, \label{eq:form-row}\\
  \textstyle\sum_{i} x_{ij} &= 1, & j &= 0,\dots,N-1, \label{eq:form-col}\\
  \textstyle\sum_{(i,j)\in D_d} x_{ij} + s^{-}_{d} &= 1, & d &= 0,\dots,2N-4, \label{eq:form-diag}\\
  \textstyle\sum_{(i,j)\in A_d} x_{ij} + s^{+}_{d} &= 1, & d &= 0,\dots,2N-4, \label{eq:form-anti}\\
  x_{ij} &= 1, & (i,j) &\in G, \label{eq:form-giv}
\end{align}
\end{subequations}
in the unknowns $x_{ij}, s^{-}_{d}, s^{+}_{d} \in \{0,1\}$. Equation~\eqref{eq:form-row} is the
representation of \secref{sec:rules} written out: each board row carries exactly one queen, which is
clauses (i) and (ii) of \defref{def:completion}. Equation~\eqref{eq:form-giv} is clause (iv). The
remaining two families are clause (iii), and they deserve an argument, because
\defref{def:attack} treats the column and the two diagonals alike whereas \eqref{eq:form} does not.

\subsection*{Why the columns admit an equality and the diagonals do not}

\begin{lemma}[the column condition is a covering condition]\label{lem:pigeon}
Let $q$ satisfy clauses (i)--(iii) of \defref{def:completion}. Then for every $j < N$ there is
exactly one $i < N$ with $q_i = j$ (\lstinline{exists_row_of_col})\qsrc{Basic.lean}{991}.
\end{lemma}

\begin{proof}
By (ii) the map $i \mapsto q_i$ sends $\{0,\dots,N-1\}$ into itself, and by (iii) with the first
disjunct of \defref{def:attack} it is injective. An injective self-map of a finite set is a
bijection, so each $j$ has exactly one preimage. Uniqueness is where the development uses it, in
\lstinline{encode_hitList_colRow}\qsrc{Basic.lean}{1045}.
\end{proof}

\lemref{lem:pigeon} is what licenses \eqref{eq:form-col}: \enquote{no two queens share a column} and
\enquote{every column carries exactly one queen} agree on boards satisfying (i)--(iii), and only the
second is a linear equation. Nothing of the sort holds for the diagonals. A board has $2N-1$
diagonals in each direction and carries only $N$ queens, so on any completion at least $N-1$ of them
are empty and the covering statement is plainly false. A diagonal constraint is therefore genuinely
\emph{at most one}, which \eqref{eq:form} cannot say directly; the standard device recovers it.

\begin{lemma}[at most one, by one auxiliary neuron]\label{lem:slack}
For a set $S$ of neurons and a fresh Boolean $s$, the constraint $\sum_{u \in S}x_u + s = 1$ is
satisfiable iff $\sum_{u \in S}x_u \le 1$, and then $s = 1 - \sum_{u \in S}x_u$.
\end{lemma}

If the diagonal is empty the auxiliary neuron takes up the slack; if it carries one queen the
auxiliary neuron is $0$; if it carries two, no value of it satisfies the row, which then contributes
at least $1$ to the objective. The development has no separate lemma for this: it is the shape of
the rows themselves, and it is used through
\lstinline{exists_diag_slack}\qsrc{Basic.lean}{268} and \lstinline{exists_anti_slack}. One auxiliary neuron suffices however long the diagonal is. This
asymmetry between \eqref{eq:form-col} and \eqref{eq:form-diag}--\eqref{eq:form-anti} is the only
point at which
\eqref{eq:form} is not a transcription of \defref{def:completion}, and every structural peculiarity
below --- the $4N-6$ surplus neurons, the non-uniform degrees, the half-integral thresholds ---
follows from it.

\subsection*{The shape of the system}

Collect~\eqref{eq:form-row}--\eqref{eq:form-giv} into a matrix $A \in \{0,1\}^{m \times n}$,
one row per equation and one column per neuron, with $A_{r,u} = 1$ exactly when the neuron $u$
occurs in the row $r$, and put $b = \mathbf{1} \in \mathbb{Z}^{m}$. Write
$\deg(u) = \sum_{r} A_{r,u}$ for the number of rows containing $u$.

\begin{proposition}[shape of the system]\label{prop:sizes}
For $N \ge 2$ the system \eqref{eq:form} has
\begin{equation}\label{eq:sizes}
  n \;=\; N^{2} + 2(2N-3) \;=\; N^{2} + 4N - 6
  \qquad\text{and}\qquad
  m \;=\; 2N + 2(2N-3) + |G| \;=\; 6N - 6 + |G| .
\end{equation}
Moreover $\deg(u) \in \{1,3,4,5\}$ for every neuron $u$, and no two distinct neurons occur together
in more than one row; equivalently, $A^{\top}\!A$ takes only the values $0$ and $1$ off its
diagonal.
\end{proposition}

\begin{proof}
The counts are immediate from \eqref{eq:form} and from there being $2N-3$ diagonals in each
direction. An auxiliary neuron occurs in the single row of its own diagonal, so its degree is $1$. A
square $(i,j)$ occurs in the row \eqref{eq:form-row} of $i$, the row \eqref{eq:form-col} of $j$, the
row \eqref{eq:form-diag} of its difference diagonal when that diagonal lies in the family
\eqref{eq:diags}, the row \eqref{eq:form-anti} of its sum diagonal likewise, and one row
\eqref{eq:form-giv} when $(i,j) \in G$. The diagonals outside \eqref{eq:diags} are those with
$i - j = \pm(N-1)$ or $i + j \in \{0, 2N-2\}$, and each consists of a single corner; the four
corners are therefore exactly the squares lying on one such diagonal, and each lies on exactly one.
So a corner has degree $3$ and every other square degree $4$, in either case one more if the square
carries a given. For the last claim: two distinct auxiliary neurons share no row; an auxiliary
neuron shares only its own diagonal row with a square; a row \eqref{eq:form-giv} contains a single
neuron; and two distinct squares lie in a common row only if they agree in board row, in column, in
difference or in sum, and any two of those four agreements force the squares to coincide.
\end{proof}

\noindent The counts (\lstinline{nvars}\qsrc{Basic.lean}{148}, \lstinline{nrows}\qsrc{Basic.lean}{155})
and the degrees
(\lstinline{rowsOf})\qsrc{Basic.lean}{320} are guarded in the development. The last claim is proved
here only, and is one of the two statements in this paper with no counterpart in it.

On the ordinary board with no givens, \eqref{eq:sizes} gives $n = 90$ and $m = 42$, and
$A^{\top}\!A$ has $852$ nonzero off-diagonal entries, $728$ joining two squares and $124$
joining a square to the auxiliary neuron of a diagonal through it. \figref{fig:inc} draws the rows
through one square.

\begin{remark}[indexing without subtraction]\label{rem:onslack}
The membership conditions of \eqref{eq:diags} are differences, and the development indexes squares by
natural numbers, on which subtraction is truncated. They are therefore tested additively
(\lstinline{onSlack})\qsrc{Basic.lean}{193}: with the $2(2N-3)$ diagonal rows numbered by $e$, the
square $(i,j)$ lies on the $e$-th exactly when
\begin{equation}\label{eq:onslack}
  i + N = e + 2 + j \quad\text{if } e < 2N-3, \qquad\qquad
  i + j + (2N-3) = e + 1 \quad\text{otherwise,}
\end{equation}
the first $2N-3$ values of $e$ indexing the $D_d$ and the rest the $A_d$. Over $\mathbb{Z}$ the two
conditions are $i - j = e + 2 - N$ and $i + j = (e - (2N-3)) + 1$, which is \eqref{eq:diags}; neither
form ever subtracts.
\end{remark}

\subsection*{The energy}

Recall from~\citeref{hnbm} that a Hopfield network on a finite neuron set $U$ with $|U| = n$ carries
a symmetric weight matrix $W = (w_{uv})_{u,v \in U}$ with zero diagonal and one threshold $\theta_u$
per neuron. A state assigns an activation to each neuron; for the $\{0,1\}$ two-state instance,
which is the one used here, we write it as a vector $x \in \{0,1\}^{n}$, and its energy is
\begin{equation}\label{eq:E}
  E(x) \;=\; -\tfrac12\sum_{u \ne v} w_{uv}\,x_u x_v \;+\; \sum_u \theta_u x_u .
\end{equation}
Write $\ell_u(x) = \sum_v w_{uv}x_v - \theta_u$ for the \emph{local field} at $u$. The network
updates one neuron at a time, setting $x_u$ to $1$ when $\ell_u(x) > 0$ and to $0$ otherwise; such
asynchronous updates never increase $E$ and reach a stable state in finitely many steps, provided
the update order is fair --- a condition that~\citeref{hnbm} found to be necessary and absent from
textbook statements. A Boltzmann machine keeps \eqref{eq:E} and makes the update stochastic: $x_u$
becomes $1$ with probability $\bigl(1 + e^{-\ell_u(x)/T}\bigr)^{-1}$ at temperature $T > 0$. At
fixed $T$ the resulting chain is ergodic, with the Boltzmann law $\propto e^{-E(x)/T}$ as its unique
stationary distribution; that uniqueness is the Perron--Frobenius result of~\citeref{hnbm} and is
what \corref{cor:gibbs} below rests on. A point of $\{0,1\}^{n}$ solves \eqref{eq:form}
exactly when the squared residual $\lVert Ax - b\rVert^2$ vanishes --- such a point is called
\emph{feasible} --- so the decision problem is the minimization of a quadratic form in binary
variables, a quadratic unconstrained binary optimization
problem~\citeref{lucas2014,glover2019}. The following identity, which holds for any
system of the shape \eqref{eq:form} and not only for the queens one, says that this form is that
energy after an affine change of scale.

\begin{proposition}[the objective is an energy]\label{prop:energy}
Let
\begin{equation}\label{eq:p}
  p(x) \;=\; \lVert Ax - b\rVert^2
        \;=\; \sum_{r=1}^{m}\Bigl(\ \sum_{u \in r} x_u - b_r\ \Bigr)^{\!2} ,
\end{equation}
and put
\begin{equation}\label{eq:WT}
  w_{uv} = -(A^{\top}\!A)_{uv} \ \ (u \ne v), \qquad
  w_{uu} = 0, \qquad
  \theta_u = \tfrac12\deg(u) - \sum_{r \ni u} b_r .
\end{equation}
Then $E(x) = \tfrac12\bigl(p(x) - \lVert b\rVert^2\bigr)$ for every $x \in \{0,1\}^{n}$
(\lstinline{zeroOneHamiltonian_eq})\qusrc{Net.lean}{353}. Hence $E$
attains its minimum $-\tfrac12\lVert b\rVert^2$ exactly at the feasible states
(\lstinline{minEnergy}, \lstinline{energy_eq_min_iff}\qusrc{Minimizers.lean}{130}), and every other
state lies at least $\tfrac12$ above it
(\lstinline{half_le_energy_sub_min}\qusrc{Minimizers.lean}{145}).
\end{proposition}

\begin{proof}
Since $A$ is $0/1$ we have $(A^{\top}\!A)_{uu} = \deg(u)$, and
$b^{\top}\!Ax = \sum_u\bigl(\sum_{r \ni u} b_r\bigr)x_u$. Expanding \eqref{eq:p},
\[
  p(x) \;=\; x^{\top}\!A^{\top}\!Ax - 2\,b^{\top}\!Ax + \lVert b\rVert^2
       \;=\; \sum_{u \ne v}(A^{\top}\!A)_{uv}x_ux_v \;+\; \sum_u \deg(u)\,x_u^{2}
             \;-\; 2\sum_u\Bigl(\sum_{r \ni u}b_r\Bigr)x_u \;+\; \lVert b\rVert^2 .
\]
The diagonal term is quadratic, and the quadratic part of \eqref{eq:E} ranges only over $u \ne v$;
since $x_u \in \{0,1\}$ we may replace $x_u^{2}$ by $x_u$, making it linear. Substituting
\eqref{eq:WT} into the three remaining sums,
\[
  \sum_{u \ne v}(A^{\top}\!A)_{uv}x_ux_v = -\sum_{u \ne v}w_{uv}x_ux_v ,
  \qquad
  \sum_u \deg(u)\,x_u - 2\sum_u\Bigl(\sum_{r \ni u}b_r\Bigr)x_u = 2\sum_u \theta_u x_u ,
\]
the second because $2\theta_u = \deg(u) - 2\sum_{r \ni u}b_r$. Hence
$p(x) = 2E(x) + \lVert b\rVert^2$. For the rest, $p$ is a sum of squares of integers, so
$p(x) \ge 0$ with equality exactly when every row of \eqref{eq:form} holds, and $p(x) \ne 0$ forces
$p(x) \ge 1$ (\lstinline{one_le_penaltyR_of_ne_zero}\qusrc{Minimizers.lean}{84}).
\end{proof}

Here $b = \mathbf{1}$, so $\lVert b\rVert^2 = m$ and $\sum_{r \ni u}b_r = \deg(u)$, whence
$\theta_u = -\tfrac12\deg(u)$ and the minimum is $-\tfrac12 m$. \eqref{eq:E} is not an analogy:
$E$ is the $\{0,1\}$ Boltzmann Hamiltonian of the formalization~\citeref{hnbm} evaluated at the
parameters \eqref{eq:WT}, so every statement that formalization proves about minimizing that
Hamiltonian is a statement about $p$ --- and, by \thmref{thm:iff} below, about the chessboard.

\begin{remark}[why the threshold is stored doubled]\label{rem:doubling}
By \propref{prop:sizes} the degree is odd at every auxiliary neuron, at every corner carrying no
given, and at every interior square carrying one, so $\theta_u = -\tfrac12\deg(u)$ is a
half-integer at all of those. Storing $\theta$ directly in an integer array
would silently restrict the encoding to even degrees; the development stores $2\theta_u = -\deg(u)$
instead (\lstinline{problem_theta})\qsrc{Basic.lean}{475} and halves it once, where the real-valued
energy is defined. The inner loop of the search
consequently computes twice the local field $Wx - \theta$. This is invisible to the deterministic
dynamics of \secref{sec:search} and not to the stochastic one; see \remref{rem:temperature}.
\end{remark}

\begin{figure}[H]
\centering
\begin{tabular}{@{}c@{\hspace{2.6em}}c@{}}
\begin{qcpic}{0.62cm}
  \qcboard{4}
  \qcrule{cRow}{1}{0}{1}{3}
  \qcrule{cCol}{0}{2}{3}{2}
  \qcrule{cDiag}{0}{1}{2}{3}
  \qcrule{cAnti}{0}{3}{3}{0}
  \qcmark{1}{2}
  \node[cRow,  font=\scriptsize] at (-0.45,-1.5) {$r_{1}$};
  \node[cDiag, font=\scriptsize] at (1.5,0.34)   {$r_{9}$};
  \node[cCol,  font=\scriptsize] at (2.5,0.34)   {$r_{6}$};
  \node[cAnti, font=\scriptsize] at (3.5,0.34)   {$r_{15}$};
\end{qcpic}
&
\begin{qcpic}{0.62cm}
  \qcboard{4}
  \qcrule{cRow}{0}{0}{0}{3}
  \qcrule{cCol}{0}{0}{3}{0}
  \qcrule{cDiag}{0}{0}{3}{3}
  \qcmark{0}{0}
  \draw[cAnti,densely dotted,line width=0.9pt] (0.5,-0.5) circle (0.36);
  \node[cRow,  font=\scriptsize] at (-0.45,-0.5) {$r_{0}$};
  \node[cCol,  font=\scriptsize] at (0.5,0.34)   {$r_{4}$};
  \node[cDiag, font=\scriptsize] at (3.5,-4.45)  {$r_{10}$};
  \node[cAnti, font=\scriptsize] at (-0.7,0.5)   {no row};
\end{qcpic}
\\[0.5em]
\qclab{(a) an interior square: $\deg = 4$} & \qclab{(b) a corner: $\deg = 3$}
\end{tabular}
\caption{The rows of $A$ through one square of the $4 \times 4$ board, on which $n = 26$ and
$m = 18$. Each capsule is one equation of~\eqref{eq:form}, drawn through the squares whose neurons
it contains: \textcolor{cRow}{board row}, \textcolor{cCol}{column}, \textcolor{cDiag}{diagonal},
\textcolor{cAnti}{anti-diagonal}. (a) The square $(1,2)$ lies in $r_1$, $r_6$, $r_9$ and $r_{15}$.
(b) The anti-diagonal through the corner $(0,0)$ is the single square $\{(0,0)\}$, drawn dotted,
which generates no equation, so the corner lies in three rows only and its threshold is
$-\tfrac32$.}\label{fig:inc}
\end{figure}

\section{What the translation is proved to do}\label{sec:proved}

The translation of \secref{sec:qubo} is a definition, and by itself it establishes nothing: a
network built from a problem by hand is only as good as the hand. If the construction is wrong then
a correct descent returns a wrong answer, and a descent that finds nothing says nothing at all. The
statements below are therefore proved rather than argued. \propref{prop:energy} has already
identified the objective with the energy; what remains is to connect both to the chessboard.
Throughout, the instance is assumed \emph{well formed}, meaning that every given lies on the board.

Let $\mathrm{enc}$ (\lstinline{encode}\qsrc{Basic.lean}{935}) send a board $q$ to the point of
$\{0,1\}^{n}$ setting $x_{i q_i}$ for each $i < N$, with each auxiliary neuron set exactly when its
diagonal is empty, and let $\mathrm{dec}$ (\lstinline{decode}\qsrc{Basic.lean}{605}) send a point $x$
to the array whose $i$-th entry is the least $j$ with $x_{ij} = 1$, or $N$ if there is none.

\begin{theorem}[soundness and completeness]\label{thm:iff}
Let $(N,G)$ be well formed. If $p(x) = 0$ then $\mathrm{dec}(x)$ is accepted by the checker
(\lstinline{decode_isQueens}\qsrc{Basic.lean}{909}); and if a board $q$ is accepted by the checker then
$p(\mathrm{enc}(q)) = 0$ (\lstinline{encode_penalty_zero}\qsrc{Basic.lean}{1179}). Consequently
\[
  \exists\, x \in \{0,1\}^{n}.\ p(x) = 0
  \qquad\Longleftrightarrow\qquad
  \exists\, q.\ q \text{ is a completion of } (N,G),
\]
the right-hand side being \defref{def:completion} by \thmref{thm:spec}
(\lstinline{exists_zero_iff_completion})\qsrc{Spec.lean}{76}.
\end{theorem}

\begin{theorem}[the minimizers]\label{thm:min}
Let $(N,G)$ be well formed. Then $E$ attains its minimum if and only if $(N,G)$ has a completion
(\lstinline{exists_minEnergy_iff})\qsrc{Converge.lean}{139},
and $\mathrm{dec}(x)$ is a completion for every $x$ attaining it
(\lstinline{isQueens_decode_of_minEnergy}\qsrc{Converge.lean}{130}).
\end{theorem}

\begin{corollary}[low-temperature concentration]\label{cor:gibbs}
Let $(N,G)$ be well formed and suppose it has a completion. Then the unique stationary distribution
of the random-scan Gibbs chain of the network \eqref{eq:WT} at temperature $T$ gives the states with
$p > 0$ a total mass that tends to $0$ as $T \to 0^{+}$
(\lstinline{stationary_infeasible_tendsto_zero})\qusrc{Gibbs.lean}{365}.
\end{corollary}

\begin{remark}[on the hypotheses]\label{rem:hyp}
Well-formedness (\lstinline{givensOk})\qsrc{Basic.lean}{626} is the only hypothesis, and it cannot be
dropped: an off-board given contributes a
row \eqref{eq:form-giv} containing no neuron at all, which no state satisfies, so $p$ would have no
zero on an instance that may well have a completion. No lower bound on $N$ is assumed by
\thmref{thm:iff} or \thmref{thm:min} --- at $N = 0$ there are no neurons and no rows, $p$ is an
empty sum, and the empty board is vacuously a completion --- although \eqref{eq:sizes} needs
$N \ge 2$. The hypothesis of \corref{cor:gibbs} is likewise not decoration: on a blocked instance no
state attains $-\tfrac12 m$, and there is nothing for the mass to concentrate on. Note also that
\corref{cor:gibbs} is a statement about the equilibrium distribution at a fixed temperature, not
about any cooling schedule, and in particular not about the annealed run of \secref{sec:search}.
\end{remark}

\thmref{thm:iff} is stated against the checker and turned into a statement about chess by
\thmref{thm:spec}, which is why that theorem was worth proving. The gap of $\tfrac12$ in
\propref{prop:energy} is what makes the objective a decision procedure rather than an
approximation: no state has energy merely close to the minimum, so a run ending above the minimum
ends at least $\tfrac12$ above it, and there is nothing to be gained by reporting how close it came.

\begin{remark}[the second half of \thmref{thm:min} does not reverse]\label{rem:slack}
The auxiliary neurons are the reason. Decoding reads only the square neurons, so a state whose board
decodes to a completion need not attain the minimum: on the empty $4 \times 4$ board, flipping one
auxiliary bit of $\mathrm{enc}(1,3,0,2)$ leaves the decoded board unchanged while raising $p$ from
$0$ to $1$. What \thmref{thm:iff} and \thmref{thm:min} give is an equivalence between the existence
of a minimizer and the existence of a completion, together with soundness of each individual
minimizer. That is what a search needs, since a search reports the state it stopped at. Like the
last clause of \propref{prop:sizes}, this one-line computation is not in the development.
\end{remark}

\subsection*{Boards with no completion}

Because \thmref{thm:iff} is an equivalence, the absence of a zero of $p$ is a statement about the
board rather than about the search. On six of the seven blocked instances of \tabref{tab:blocked}
that statement is proved, by exhausting the tuples the givens leave open using kernel reduction;
\lstinline{native_decide} is not used, so no compiler is part of the trusted base.

\figref{fig:hard} shows the two extremes. \texttt{Attack} needs no enumeration at all: its two
givens share a diagonal, so the checker's diagonal clause refutes the board outright
(\lstinline{attacking_no_queens}\qsrc{Examples.lean}{192}). \texttt{Blk7}
is the instance worth noting. Its givens share no row, no column and no diagonal, so the placement
is locally unobjectionable and no attack is displayed, and the board admits no completion regardless;
this is the phenomenon that makes Completion hard rather than merely large. A check over
$7^{5} = 16807$ tuples settles it (\lstinline{blk7_no_queens}\qsrc{Bench.lean}{388}), and by
\thmref{thm:iff} the refutation extends to every one of the $2^{71}$ states of that instance
(\lstinline{blk7_no_zero}\qsrc{Bench.lean}{403}). The seventh, \texttt{Blk8}, leaves $8^{6} = 262144$ tuples
open, which is beyond what kernel reduction will do here; it is reported as an observation of the
baseline rather than as a theorem, and \tabref{tab:blocked} marks it as such.

\begin{figure}[H]
\centering
\begin{tabular}{@{}c@{\hspace{2.6em}}c@{}}
\begin{qcpic}{0.46cm}
  \qcboard{8}\qcattack{0}{0}{1}{1}\qcgivens{0/0,1/1}
\end{qcpic}
&
\begin{qcpic}{0.46cm}
  \qcboard{7}\qcgivens{0/0,1/6}
\end{qcpic}
\\[0.5em]
\qclab{(a) \texttt{Attack}} & \qclab{(b) \texttt{Blk7}}
\end{tabular}
\caption{Two instances with no completion. (a) The two givens attack each other along a diagonal,
shown the way the board widget shows it, and the checker refutes the board with no enumeration.
(b) The givens of \texttt{Blk7} share no row, no column and no diagonal, so nothing is displayed and
nothing is locally wrong; the board admits no completion all the same.}\label{fig:hard}
\end{figure}

\subsection*{Where the certification stops}

The objective is identified with the energy by \propref{prop:energy}; the integer routine that computes
the net input in the inner loop of the search is proved equal to twice the network's local field
(\lstinline{netVec_eq_localField}\qusrc{Refine.lean}{256}), the factor of two being the one that lets
$\theta$ be stored doubled; and the minimizers are
characterized by \thmref{thm:min}. The \emph{trajectory} is not certified, and it departs from the
certified dynamics in two independent ways. The recurrences of \secref{sec:search} accumulate their
net input over time, and the accumulated $u$ is not the local field of the current state --- the two
coincide only at $t = 0$; and the update is synchronous, whereas every convergence, Lyapunov and
detailed-balance result of~\citeref{hnbm} is stated for the single-site update. The second gap can
be closed on its own: the development implements an asynchronous variant, which is the one the
theory covers, at the cost of a slower cooling schedule and roughly three times the wall clock. The
first cannot, so even that variant would not be the dynamics of~\citeref{hnbm}. Convergence of
\eqref{eq:dyn} to local optima is asserted in~\citeref{liwang2022} on the strength of simulation,
and remains unproved here.

What a run of \secref{sec:search} establishes when it succeeds is therefore a theorem about the
board; what it establishes when it fails is nothing, and the negative results above come from
enumeration instead.


\section{The collaborative search}\label{sec:search}

A single Hopfield network stops at the first state that no single update improves, which on these
instances is almost never a minimum. Collaborative neurodynamic
optimization~\citeref{che2019} improves on it by running a population of networks concurrently and
re-initializing their states by a particle swarm rule~\citeref{kennedy1995}, and it is the
scheme of~\citeref{liwang2022} that we apply here, unchanged in structure.

Two recurrences are used, both driven by the same accumulated net input. Starting from a seed
$x(0)$ and $u(0) = 0$,
\begin{subequations}\label{eq:dyn}
\begin{align}
  u(t+1) &= u(t) + Wx(t) - \theta, \label{eq:dyn-u}\\
  x_i(t+1) &= \begin{cases} 1 & \text{if } u_i(t+1) > 0,\\ 0 & \text{otherwise,}\end{cases}
    \label{eq:dyn-dhnm}\\
  \Pr\bigl[x_i(t+1) = 1\bigr] &= \bigl(1 + e^{-u_i(t+1)/T(t)}\bigr)^{-1}, \label{eq:dyn-bmm}
\end{align}
\end{subequations}
with $T(t) = T_0\eta^{t}$ a geometric cooling schedule over the sweeps of one run. A discrete
Hopfield network with a momentum term (DHNm) is \eqref{eq:dyn-u} with \eqref{eq:dyn-dhnm}; a
Boltzmann machine with one (BMm) is \eqref{eq:dyn-u} with \eqref{eq:dyn-bmm}. These are equations
(3) and (6) of~\citeref{liwang2022}, whose equation (2) is the threshold \eqref{eq:dyn-dhnm}; in the
development they are \lstinline{dhnmStep}\qusrc{Dynamics.lean}{179} and
\lstinline{bmmStep}\qusrc{Dynamics.lean}{213}. Note
that $u(t)$ is an
accumulated field and not the local field of \secref{sec:qubo}: the two agree only at $t = 0$, where
$u(1) = Wx(0) - \theta$. \figref{fig:alg} states the search around them (\lstinline{search}\qusrc{Search.lean}{171}). The only
problem-specific input is the list of mutually exclusive neuron groups used at step 1, which here is
the $N$ squares of each board row --- the rows \eqref{eq:form-row} of \figref{fig:inc}.

\begin{remark}[three properties of \eqref{eq:dyn} worth stating]\label{rem:temperature}
First, the update is \emph{synchronous}: every neuron is recomputed from $u(t+1)$ at once.
\citeref{liwang2022} calls its models \enquote{activated synchronously}, but that phrase is about
the population rather than the neurons within one network, so the reading is ours; the momentum
term of \eqref{eq:dyn-u} is the classical device for taming the oscillation a synchronous update
invites, which is the reading it supports. An asynchronous single-site variant is also
implemented, and \secref{sec:proved} returns to what it would buy.

Second, \eqref{eq:dyn-u} is an \emph{undamped} accumulator, so the recurrence need not converge and
a synchronous update can fall into a $2$-cycle. Each inner run therefore returns the state it
happens to end on after its fixed number of sweeps; the implementation can instead track the best
state visited, but the runs of \secref{sec:search} do not, so the value reported for a copy is the
value of its last state and not a minimum over its trajectory.

Third, by \remref{rem:doubling} the inner loop computes $2(Wx - \theta)$, so the accumulator it
carries is $2u$. The threshold test \eqref{eq:dyn-dhnm} is unaffected, since $2z > 0$ iff $z > 0$;
the logistic of \eqref{eq:dyn-bmm} is not, and the implementation compensates by building its table
at $2T_0$ (\lstinline{ModelConfig.tabulate}\qusrc{Dynamics.lean}{169}), so that the law it samples is
exactly \eqref{eq:dyn-bmm} at $T(t)$. Getting this wrong anneals from $T_0/2$, which is why the
development guards the two conventions against each other.
\end{remark}

\begin{figure}[H]
\centering\small
\setlength{\tabcolsep}{0pt}
\begin{tabular}{@{}p{0.965\textwidth}@{}}
\toprule
\textbf{The collaborative search} (Algorithm 2 of~\citeref{liwang2022}), as applied
to~\eqref{eq:form}.\\
\midrule
\textbf{Input:} the network $(W,\theta)$ of~\eqref{eq:WT}; a model $\mathcal{M} \in
\{\text{DHNm},\ \text{BMm}\}$; and $K$, $M$, $c_0$, $c_1$, $c_2$, $P_m$, $\mathcal{T}$, $T_0$,
$\eta$.\\
\textbf{Output:} the incumbent $x^{*}$ and its objective value $p(x^{*})$.\\
\midrule
\emph{1} \quad draw $x^{(i)}(0)$ with one neuron set in each board row, $v^{(i)} \sim U[-1,1]^{n}$;
set $\hat x^{(i)} \leftarrow x^{(i)}(0)$ and $x^{*} \leftarrow \arg\min_i p(\hat x^{(i)})$\\
\emph{2--7} \quad\, \textbf{repeat:} run each of the $K$ copies of $\mathcal{M}$ from $x^{(i)}(0)$
for a fixed number of sweeps of~\eqref{eq:dyn}, giving $x^{(i)}$; set
$\hat x^{(i)} \leftarrow x^{(i)}$ whenever $p(x^{(i)}) < p(\hat x^{(i)})$\\
\emph{8--14} \quad $x^{*} \leftarrow \arg\min_i p(\hat x^{(i)})$ if that improves on $x^{*}$, and
reset the stall counter; otherwise increment it. Exit if $p(x^{*}) = 0$\\
\emph{15--20} \, $v^{(i)} \leftarrow c_0 v^{(i)} + c_1 r_1 (\hat x^{(i)} - x^{(i)})
+ c_2 r_2 (x^{*} - x^{(i)})$ with $r_1, r_2 \sim U[0,1]$, then
$x^{(i)}(0) \leftarrow \mathrm{round}\bigl(x^{(i)}(0) + v^{(i)}\bigr)$\\
\emph{21--24} \, $\delta \leftarrow \frac{1}{Kn}\sum_i \lVert \hat x^{(i)} - x^{*}\rVert_2$;
\ \textbf{if} $\delta\sqrt{n} < \mathcal{T}$ \textbf{then} flip each bit of every $x^{(i)}(0)$ with
probability $P_m$\\
\quad\ \ \ \ \ \ \ \, \textbf{until} $p(x^{*}) = 0$, or the stall counter exceeds $M$, or the round
cap is reached\\
\bottomrule
\end{tabular}
\caption{The search, with the step numbers of~\citeref{liwang2022}; steps 15--20 are its equation
(7), the diversity of step 21 its equation (8), and the mutation of step 23 its equation (9). Four
points a literal reading does not settle are resolved here and flagged rather than hidden: the
one-hot draw at step 1; the $\sqrt{n}$ at step 22; the target of the mutation at step 23; and the
reading of step 16, whose equation (7) we take with the current position as the just-computed
equilibrium $x^{(i)}$ and with $\hat x^{(i)}$ as the personal best, so that both the cognitive and
the social term are measured from the equilibrium. All four are discussed below. Note that (8) uses
$\lVert\cdot\rVert_2$ unsquared, so each term is the square root of a Hamming
distance.}\label{fig:alg}
\end{figure}

\subsection*{Parameters}

Running the search requires a numerical value for every parameter in \tabref{tab:params}. For eight
of them --- $T_0$, $\eta$, the inner-iteration count, $P_m$, $\mathcal{T}$ and the swarm constants
$c_0$, $c_1$, $c_2$ --- \citeref{liwang2022} gives no numerical value anywhere, so all eight were
recovered experimentally; the population size, the stall limit and the round cap are the settings
used here. Two of the recovered eight are not merely a matter of tuning.

The social coefficient $c_2$ weights the pull of the incumbent. If it is comparable to $c_1$ then
every copy is drawn onto $x^{*}$, and once $x^{*}$ lies on a plateau --- which on these instances
happens almost at once --- the population stops exploring, and further copies are further copies of
the same stuck point. Taking $c_2$ well below $c_1$ removes this. The diversity threshold
$\mathcal{T}$ has to be made scale-free: the diversity measure of~\citeref{liwang2022} is an average
of Euclidean distances divided by the number of neurons, hence at most $1/\sqrt{n}$, so a fixed
threshold means something different on every instance and on a large one fires every round. We
compare it against the diversity scaled by $\sqrt{n}$, as step 22 of \figref{fig:alg} records. The
recovered value is model-specific, and it transfers: on the empty $8 \times 8$ board DHNm succeeds
on four of ten seeds at BMm's $0.4$ against nine of ten at its own $0.9$.

Three further choices are interpretations rather than numbers. Step 1 of~\citeref{liwang2022} says
only that the initial state is binary; uniform bits set about half the neurons, whereas a feasible
point sets exactly one per board row, so a uniformly seeded run spends its whole equilibration
merely reaching one queen per row. Seeding one-hot starts every copy on the feasible set
of~\eqref{eq:form-row} and leaves the dynamics the columns and the diagonals. Step 23 says only to
mutate, without naming the target; we mutate the seeds $x^{(i)}(0)$, never $x^{*}$, since mutating
the incumbent would forfeit the best board found. And step 16 writes its two difference terms
against \enquote{the $i$-th particle}, which the surrounding steps use for both the equilibrium
state and the personal best; we read the base point as the equilibrium and the cognitive target as
the personal best, as \figref{fig:alg} records.

\begin{table}[H]
\centering\small
\setlength{\tabcolsep}{5pt}
\begin{tabular}{@{}l p{0.44\textwidth} l@{}}
\toprule
& meaning & value\\
\midrule
$K$            & population size                                     & $20$\\
$M$            & stall limit, in rounds                              & $20$\\
               & round cap                                           & $60$\\
               & inner iterations per equilibration                  & $40$\\
$T_0$          & initial temperature                                 & $3.0$\\
$\eta$         & cooling rate                                        & $0.9$\\
$c_0$          & inertia weight                                      & $0.5$\\
$c_1$          & cognitive coefficient                               & $2.0$\\
$c_2$          & social coefficient                                  & $0.25$\\
$P_m$          & bit-flip mutation probability                       & $0.05$\\
$\mathcal{T}$  & diversity threshold, as a fraction of $\delta_{\max} = 1/\sqrt{n}$
               & $0.4$ (BMm), $0.9$ (DHNm)\\
\bottomrule
\end{tabular}
\caption{The parameters of \figref{fig:alg}, used for every collaborative run in
\tabref{tab:results} and \tabref{tab:blocked}, and collected in
\texttt{SearchConfig}\qusrc{Search.lean}{61}. The eight named in \secref{sec:search} are given no
numerical value anywhere in~\citeref{liwang2022} and were recovered
experimentally.}\label{tab:params}
\end{table}

\subsection*{Results}

The parameters are fixed across every instance and so is the base seed. The generator is
\texttt{splitmix64} and no clock is read, so both tables are reproducible exactly. The baseline is a
chronological backtracking completion search written independently of the encoding. It is itself
certified, in that every board it returns is proved to satisfy the checker
(\lstinline{solve_isQueens}\qsrc{Baseline.lean}{223}), and the column \textsf{ok} of
\tabref{tab:results} is that theorem being re-checked at run time; the converse is
not proved, so the baseline's failures would be observations, and it has none.

\begin{table}[H]
\centering\footnotesize
\setlength{\tabcolsep}{3.5pt}
\begin{tabular}{@{}lrrrr rr rrr r r@{}}
\toprule
& & & & & \multicolumn{2}{c}{backtr.}
& \multicolumn{3}{c}{CNS/BMm} & DHNm & 1 BMm\\
\cmidrule(lr){6-7}\cmidrule(lr){8-10}\cmidrule(lr){11-11}\cmidrule(lr){12-12}
board & $N$ & $|G|$ & $n$ & $m$ & ok & \#compl. & hit & rounds & $s/k$ & $s/k$ & $s/20$\\
\midrule
\texttt{Q4}     & 4  & 0 & 26  & 18 & y & 2   & y & 1  & 10/10 & 10/10 & 7\\
\texttt{Q5}     & 5  & 0 & 39  & 24 & y & 10  & y & 1  & 10/10 & 10/10 & 5\\
\texttt{Q6}     & 6  & 0 & 54  & 30 & y & 4   & y & 4  & 9/10  & 8/10  & 0\\
\texttt{Q7}     & 7  & 0 & 71  & 36 & y & 40  & y & 10 & 10/10 & 10/10 & 0\\
\texttt{Q8}     & 8  & 0 & 90  & 42 & y & 92  & y & 3  & 10/10 & 9/10  & 0\\
\texttt{Q9}     & 9  & 0 & 111 & 48 & y & 352 & y & 24 & 4/5   & 5/5   & 0\\
\texttt{Q10}    & 10 & 0 & 134 & 54 & y & 724 & n & 31 & 0/5   & 1/5   & 0\\
\midrule
\texttt{Comp6}  & 6  & 1 & 54  & 31 & y & 1   & n & 22 & 6/10  & 7/10  & 0\\
\texttt{Comp7}  & 7  & 1 & 71  & 37 & y & 4   & y & 8  & 10/10 & 6/10  & 0\\
\texttt{Comp8}  & 8  & 3 & 90  & 45 & y & 1   & y & 14 & 6/10  & 3/10  & 0\\
\texttt{Comp8b} & 8  & 5 & 90  & 47 & y & 1   & y & 4  & 10/10 & 10/10 & 0\\
\bottomrule
\end{tabular}
\caption{The eleven completable instances. Columns: board size; number of givens; the neuron and row
counts $n$ and $m$ of~\eqref{eq:sizes}; whether the checker accepts the board backtracking returned;
the total number of completions; whether the collaborative search driving BMm reached $p = 0$ at the
base seed, and in how many rounds; its successes over $k$ seeds; the same driving DHNm; and
successes over $20$ seeds for a single BMm run with no population at all. That last run is one
annealed pass of~\eqref{eq:dyn}, $400$ sweeps at $T = 3.0 \cdot 0.99^{t}$ from the same base seed,
and is the one configuration here not covered by \tabref{tab:params}. Backtracking found a
completion on every row. The table is printed by
\texttt{Bench.table}\qsrc{Bench.lean}{247}.}\label{tab:results}
\end{table}

\begin{table}[H]
\centering\footnotesize
\setlength{\tabcolsep}{4.5pt}
\begin{tabular}{@{}lrrrr rr rr@{}}
\toprule
board & $N$ & $|G|$ & $n$ & $m$ & \#compl. & refuted & tuples & best $p$\\
\midrule
\texttt{Q2}     & 2 & 0 & 6  & 6  & 0 & y & $2^2 = 4$      & 1\\
\texttt{Q3}     & 3 & 0 & 15 & 12 & 0 & y & $3^3 = 27$     & 1\\
\texttt{Blk4}   & 4 & 1 & 26 & 19 & 0 & y & $4^3 = 64$     & 1\\
\texttt{Blk6}   & 6 & 1 & 54 & 31 & 0 & y & $6^5 = 7776$   & 1\\
\texttt{Blk7}   & 7 & 2 & 71 & 38 & 0 & y & $7^5 = 16807$  & 1\\
\texttt{Attack} & 8 & 2 & 90 & 44 & 0 & y & none needed    & 1\\
\texttt{Blk8}   & 8 & 2 & 90 & 44 & 0 & n & $8^6 = 262144$ & 1\\
\bottomrule
\end{tabular}
\caption{The seven blocked instances. \textsf{refuted} is whether the absence of a completion is a
theorem of the development rather than a verdict of the baseline; \textsf{tuples} is the enumeration
that establishes it, over the rows the givens leave open. On every row the search found no zero on
any seed, and its best value was exactly $1$, the least it can be by \propref{prop:energy}. The
table is printed by \texttt{Bench.tableBlocked}\qsrc{Bench.lean}{292}.}\label{tab:blocked}
\end{table}

Four things are worth drawing out of the two tables.

\emph{1) The baseline wins.} Backtracking solves every completable instance and refutes every
blocked one, in times not worth tabulating. No speed claim is made for the neurodynamic algorithms
and nothing here should be read as one.

\emph{2) The population is not a refinement but the difference between working and not working.} A
single BMm run with no population succeeds on $7$ of $20$ seeds at $N = 4$ and $5$ of $20$ at
$N = 5$, and on nothing from $N = 6$ upwards or on any instance with givens. The collaborative
search reaches a completion on ten of the eleven rows under BMm and on all eleven under DHNm.

\emph{3) The two models separate, and only where there are givens.} On the seven empty boards DHNm
and BMm are indistinguishable, at $53$ of $60$ seeds each. On the four boards with givens BMm is
ahead, $32$ of $40$ against $26$. The separation is not uniform --- DHNm is the only configuration to
reach $N = 10$ at all, once in five seeds --- and it is problem-specific: on a graph-colouring corpus
encoded in the same framework the two are indistinguishable throughout.

\emph{4) The corpus is what bounds the claim.} $N = 10$ is the edge: BMm fails on it on every seed
and DHNm succeeds on one of five, and nothing larger was run. A purpose-built Hopfield network for
$n$-queens~\citeref{mandziuk1995} reaches $N = 80$, on the plain problem rather than on Completion.

Finally, a word on what is not reported. It is usual to give the best and worst objective values
over runs with their mean and standard deviation. That statistic is uninformative here: by
\propref{prop:energy} the objective is integer-valued with a gap of $\tfrac12$ in the energy, so a run
ends either at $p = 0$ or at $p \ge 1$, and on every blocked instance every failed run ends at
exactly $1$, as the last column of \tabref{tab:blocked} records. A mean over runs would report
little beyond the success rate, which is given directly, and a standard deviation would report the
same thing twice.

\section{Concluding remarks}\label{sec:concl}

We have presented a straightforward translation of an $n$-Queens Completion instance into the energy
of a $\{0,1\}$ Boltzmann machine, and certified it. The translation is parametric in $N$ and
polynomial in the size of the board --- $N^{2} + 4N - 6$ neurons and $6N - 6 + |G|$ rows --- and
since Completion is \textsf{NP}-complete~\citeref{gent2017}, no algorithm with better complexity is
known. Collaborative neurodynamic algorithms then apply to it without modification, and on eleven
completable instances they find a completion on all but one.

The translation can also be used to enumerate: the certified baseline counts the completions of
every instance (\lstinline{count}\qsrc{Baseline.lean}{84}), and on the empty boards it returns $1, 1, 0, 0, 2, 10, 4, 40, 92, 352, 724$ for
$N = 0,\dots,10$, the classical sequence~\citeref{oeis170}. That count is a cross-check on the baseline rather than a
theorem about it --- what is proved is that every board it returns is a completion, not that it
returns them all.

Particularly worth remarking is how little of this was implementation effort. The energy, the
network, both recurrences of~\eqref{eq:dyn}, the collaborative search and the whole minimizer theory
came from the existing formalization~\citeref{hnbm} unchanged; only the constraints~\eqref{eq:form}
had to be specified, in about a hundred and fifty lines. The proof that those constraints say what
they appear to say is the rest of the development, and it is not small. That asymmetry is the point:
what the proof buys is not speed --- backtracking outperforms the neurodynamic algorithms on every
instance, and a purpose-built network~\citeref{mandziuk1995} outperforms them at sizes this encoding
does not reach --- but the guarantee that a run which ends at the minimum has established a property
of the chessboard, and that the objective being descended is provably the energy of the network said
to be descending it.

Two directions remain. The trajectory is not certified, as \secref{sec:proved} sets out: the
recurrences~\eqref{eq:dyn} accumulate their net input and update synchronously, and so are not the
single-site dynamics for which convergence is proved. Switching to the asynchronous variant already
implemented would close the second gap but not the first; a solver restricted to single-site
updates \emph{without} the momentum term would be certified from the board to the last update, and
would be a different algorithm from the one~\citeref{liwang2022} proposes. And the variable-reduction stage of~\citeref{liwang2022},
which fixes the neurons the givens determine before the search begins, has no counterpart here: every
result above is obtained on the full $N^{2} + 4N - 6$ neurons. An analogue exists in principle,
since a row or column left with a single square unattacked by a given must carry a queen there. Until
it is implemented and certified, the results here and those of~\citeref{liwang2022} are not
comparable in either direction.

\section*{Acknowledgements}
% TODO

\begin{thebibliography}{99}

\bibitem{hopfield1982}
J.~J.~Hopfield.
\newblock Neural networks and physical systems with emergent collective computational abilities.
\newblock \emph{Proceedings of the National Academy of Sciences}, 79(8):2554--2558, 1982.

\bibitem{ackley1985}
D.~H.~Ackley, G.~E.~Hinton, and T.~J.~Sejnowski.
\newblock A learning algorithm for Boltzmann machines.
\newblock \emph{Cognitive Science}, 9(1):147--169, 1985.

\bibitem{liwang2022}
H.~Li and J.~Wang.
\newblock Collaborative neurodynamic algorithms for solving Sudoku puzzles.
\newblock In \emph{12th International Conference on Information Science and Technology (ICIST)},
pages 8--17. IEEE, 2022.

\bibitem{bell2009}
J.~Bell and B.~Stevens.
\newblock A survey of known results and research areas for $n$-queens.
\newblock \emph{Discrete Mathematics}, 309(1):1--31, 2009.

\bibitem{kennedy1995}
J.~Kennedy and R.~Eberhart.
\newblock Particle swarm optimization.
\newblock In \emph{Proceedings of the IEEE International Conference on Neural Networks}, volume 4,
pages 1942--1948, 1995.

\bibitem{che2019}
H.~Che and J.~Wang.
\newblock A collaborative neurodynamic approach to global and combinatorial optimization.
\newblock \emph{Neural Networks}, 114:15--27, 2019.

\bibitem{hnbm}
M.~Cipollina, M.~Karatarakis, and F.~Wiedijk.
\newblock Formalized Hopfield networks and Boltzmann machines.
\newblock In \emph{LPAR-26}, EPiC Series in Computing, 2026. To appear. arXiv:2512.07766.

\bibitem{lean4}
L.~de~Moura and S.~Ullrich.
\newblock The Lean 4 theorem prover and programming language.
\newblock In \emph{CADE-28}, volume 12699 of \emph{Lecture Notes in Computer Science}, pages
625--635. Springer, 2021.

\bibitem{gent2017}
I.~P.~Gent, C.~Jefferson, and P.~Nightingale.
\newblock Complexity of $n$-Queens Completion.
\newblock \emph{Journal of Artificial Intelligence Research}, 59:815--848, 2017.

\bibitem{mandziuk1995}
J.~Ma\'{n}dziuk.
\newblock Solving the $N$-Queens problem with a binary Hopfield-type network.
\newblock \emph{Biological Cybernetics}, 72:439--445, 1995.

\bibitem{lucas2014}
A.~Lucas.
\newblock Ising formulations of many NP problems.
\newblock \emph{Frontiers in Physics}, 2:5, 2014.

\bibitem{glover2019}
F.~Glover, G.~Kochenberger, and Y.~Du.
\newblock Quantum bridge analytics I: a tutorial on formulating and using QUBO models.
\newblock arXiv:1811.11538, 2019.

\bibitem{oeis170}
OEIS Foundation Inc.
\newblock Sequence A000170: number of ways of placing $n$ non-attacking queens on an
$n \times n$ board.
\newblock \emph{The On-Line Encyclopedia of Integer Sequences}, 2026.
\url{https://oeis.org/A000170}.

\end{thebibliography}

\end{document}
